
\documentclass{article}

\usepackage{fancyhdr}
\usepackage{extramarks}
\usepackage{amsmath}
\usepackage{amsthm}
\usepackage{amsfonts}
\usepackage{tikz}
\usepackage[plain]{algorithm}
\usepackage{algpseudocode}
\usepackage{amsmath,amssymb}

\usetikzlibrary{automata,positioning}

%
% Basic Document Settings
%

\topmargin=-0.45in
\evensidemargin=0in
\oddsidemargin=0in
\textwidth=6.5in
\textheight=9.0in
\headsep=0.25in

\linespread{1.1}

\pagestyle{fancy}
\lhead{\hmwkAuthorName}
\chead{\hmwkClass: \hmwkTitle}
\rhead{}
\lfoot{\lastxmark}
\cfoot{\thepage}

\renewcommand\headrulewidth{0.4pt}
\renewcommand\footrulewidth{0.4pt}

\setlength\parindent{0pt}

%
% Create Problem Sections
%

\newcommand{\enterProblemHeader}[1]{
    \nobreak\extramarks{}{Problem \arabic{#1} continued on next page\ldots}\nobreak{}
    \nobreak\extramarks{Problem \arabic{#1} (continued)}{Problem \arabic{#1} continued on next page\ldots}\nobreak{}
}

\newcommand{\exitProblemHeader}[1]{
    \nobreak\extramarks{Problem \arabic{#1} (continued)}{Problem \arabic{#1} continued on next page\ldots}\nobreak{}
    \stepcounter{#1}
    \nobreak\extramarks{Problem \arabic{#1}}{}\nobreak{}
}

\setcounter{secnumdepth}{0}
\newcounter{partCounter}
\newcounter{homeworkProblemCounter}
\setcounter{homeworkProblemCounter}{1}
\nobreak\extramarks{Problem \arabic{homeworkProblemCounter}}{}\nobreak{}



% Define theorem environment
\newtheorem*{theorem}{Theorem}

%
% Homework Problem Environment
%
% This environment takes an optional argument. When given, it will adjust the
% problem counter. This is useful for when the problems given for your
% assignment aren't sequential. See the last 3 problems of this template for an
% example.
%
\newenvironment{homeworkProblem}[1][-1]{
    \ifnum#1>0
        \setcounter{homeworkProblemCounter}{#1}
    \fi
    \section{Problem \arabic{homeworkProblemCounter}}
    \setcounter{partCounter}{1}
    \enterProblemHeader{homeworkProblemCounter}
}{
    \exitProblemHeader{homeworkProblemCounter}
}

%
% Homework Details
%   - Title
%   - Due date
%   - Class
%   - Section/Time
%   - Instructor
%   - Author
%


\newcommand{\hmwkTitle}{Problem set\ 10}
\newcommand{\hmwkDueDate}{November 17, 2025}
\newcommand{\hmwkClass}{Number Theory}
\newcommand{\hmwkClassTime}{Section 2}
\newcommand{\hmwkClassInstructor}{Dr. Eleanor McSpirit}
\newcommand{\hmwkAuthorName}{\textbf{Joseph Quinn}}

%
% Title Page
%

\title{
    \vspace{2in}
    \textmd{\textbf{\hmwkClass:\ \hmwkTitle}}\\
    \normalsize\vspace{0.1in}\small{\hmwkDueDate}\\
    \vspace{0.1in}\large{\textit{\hmwkClassInstructor\ \hmwkClassTime}}
    \vspace{3in}
}

\author{\hmwkAuthorName}
\date{}

\renewcommand{\part}[1]{\textbf{\large Part \Alph{partCounter}}\stepcounter{partCounter}\\}

%
% Various Helper Commands
%

% Useful for algorithms
\newcommand{\alg}[1]{\textsc{\bfseries \footnotesize #1}}

% For derivatives
\newcommand{\deriv}[1]{\frac{\mathrm{d}}{\mathrm{d}x} (#1)}

% For partial derivatives
\newcommand{\pderiv}[2]{\frac{\partial}{\partial #1} (#2)}

% Integral dx
\newcommand{\dx}{\mathrm{d}x}

% Alias for the Solution section header
\newcommand{\solution}{\textbf{\large Solution}}

% Probability commands: Expectation, Variance, Covariance, Bias
\newcommand{\E}{\mathrm{E}}
\newcommand{\Var}{\mathrm{Var}}
\newcommand{\Cov}{\mathrm{Cov}}
\newcommand{\Bias}{\mathrm{Bias}}

%  proof-step macro:
\newcommand{\step}[2]{& #1 & & \text{#2} \\}

\begin{document}

\maketitle
\pagebreak


\begin{homeworkProblem}
	Since $p = 23$, we have
	\[
		\frac{p-1}{2} = 11.
	\]
	We consider the multiples
	\[
		12,\; 24,\; 36,\; 48,\; 60,\; 72,\; 84,\; 96,\; 108,\; 120,\; 132
	\]
	and reduce each into the symmetric residue system
	\[
		\{-11,-10,\dots,-1,0,1,\dots,10,11\}.
	\]

	\[
		\begin{aligned}
			12  & \equiv -11 \pmod{23}, \\
			24  & \equiv 1 \pmod{23},   \\
			36  & \equiv -10 \pmod{23}, \\
			48  & \equiv 2 \pmod{23},   \\
			60  & \equiv -9 \pmod{23},  \\
			72  & \equiv 3 \pmod{23},   \\
			84  & \equiv -8 \pmod{23},  \\
			96  & \equiv 4 \pmod{23},   \\
			108 & \equiv -7 \pmod{23},  \\
			120 & \equiv 5 \pmod{23},   \\
			132 & \equiv -6 \pmod{23}.
		\end{aligned}
	\]

	Among these eleven residues, the negative ones are
	\[
		-11,\; -10,\; -9,\; -8,\; -7,\; -6,
	\]
	so $g = 6$. By Gauss' Lemma,
	\[
		\left( \frac{12}{23} \right) = (-1)^g = (-1)^6 = 1.
	\]

	Therefore, $12$ is a quadratic residue modulo $23$.
\end{homeworkProblem}

\begin{homeworkProblem}
	We will use Euler's Criterion:
	\[
		\left( \frac{a}{p} \right) \equiv a^{\frac{p-1}{2}} \pmod{p}.
	\]


	Since $p = 97$, Euler's Criterion requires computing $1248^{48} \pmod{97}$.

	First reduce $1248$ modulo $97$:
	\[
		1248 \equiv 84 \pmod{97}.
	\]

	And simplifying gives us \[
		84^{48} \equiv \pmod{97}
	\]

	\[
		\begin{aligned}
			84^2    & = 7056 \equiv 72 \pmod{97},             \\[4pt]
			84^4    & \equiv 72^2 = 5184 \equiv 43 \pmod{97}, \\[4pt]
			84^8    & \equiv 43^2 = 1849 \equiv 6 \pmod{97},  \\[4pt]
			84^{16} & \equiv 6^2 = 36 \pmod{97},              \\[4pt]
			84^{32} & \equiv 36^2 = 1296 \equiv 35 \pmod{97}.
		\end{aligned}
	\]

	Now I will combine the powers to get:
	\[
		84^{48} = 84^{32} \cdot 84^{16}
		\equiv 35 \cdot 36 = 1260 \pmod{97}.
	\]
	Reduce:
	\[
		1260 \equiv 96 \equiv -1 \pmod{97}.
	\]

	Therefore, $1248$ is a quadratic non-residue modulo $97$.
\end{homeworkProblem}

\begin{homeworkProblem}
	\begin{proof}

		By Gauss' Lemma,
		\[
			\left( \frac{a}{p} \right) = (-1)^{g_p},
			\qquad
			\left( \frac{a}{q} \right) = (-1)^{g_q},
		\]
		where $g_p$ is the number of negative representatives of
		\[
			a, 2a, \dots, \frac{p-1}{2}a
		\]
		in the symmetric residue system.

		Because $p = q + 4k$, the bounds of the symmetric intervals satisfy
		\[
			\frac{p-1}{2} = \frac{q-1}{2} + 2k.
		\]
		The extra $2k$ multiples of $a$ contribute values of the form
		\[
			\left(\frac{q+1}{2}a\right),
			\left(\frac{q+3}{2}a\right), \dots,
			\left(\frac{q-1}{2}+2k\right)a.
		\]

		When reduced modulo $p$ or modulo $q$, each of these additional contribute an even number of 
		negative representatives.  Therefore:
		\[
			g_p \equiv g_q \pmod{2}.
		\]
		Thus
		\[
			(-1)^{g_p} = (-1)^{g_q},
		\]
		and 
		\[
			\left( \frac{a}{p} \right)
			=
			\left( \frac{a}{q} \right).
		\]

	\end{proof}
\end{homeworkProblem}

\begin{homeworkProblem}

	\[
		2 = 1^{2} + 1^{2},
	\]
	\[
		10 = 1^{2} + 3^{2},
	\]
	\[
		18 = 3^{2} + 3^{2},
	\]
	\[
		20 = 2^{2} + 4^{2}.
	\]
\end{homeworkProblem}

\begin{homeworkProblem}
\begin{proof}
Since $p \mid (ax^{2} + bxy + cy^{2})$ and $p \nmid a$, we complete the square:
\[
ax^{2} + bxy + cy^{2}
= a\left(x + \frac{b}{2a}y\right)^{2}
 - a\left(\frac{b^{2}-4ac}{4a^{2}}\right)y^{2}.
\]
Reducing modulo $p$ gives
\[
a\left(x + \frac{b}{2a}y\right)^{2}
 \equiv \frac{D}{4a}\, y^{2} \pmod p.
\]
Multiplying both sides by $4a$, which is invertible mod $p$, yields
\[
(2ax + by)^{2} \equiv D\,y^{2} \pmod p.
\]
Define
\[
X := 2ax + by.
\]
Then we have the key congruence
\[
X^{2} \equiv D\,y^{2} \pmod p.
\]

Because $p \mid (ax^{2} + bxy + cy^{2})$ and $p \nmid a$, the congruence above implies
$p \nmid y$ (otherwise $p \mid x$, contradicting $(x,y)=1$). Hence $X \not\equiv 0 \pmod p$
and $y \not\equiv 0 \pmod p$.

Taking Legendre symbols,
\[
\left( \frac{X^{2}}{p} \right)
=
\left( \frac{D\,y^{2}}{p} \right).
\]
Using multiplicativity of the Legendre symbol,
\[
\left( \frac{X}{p} \right)^{2}
=
\left( \frac{D}{p} \right)
\left( \frac{y}{p} \right)^{2}.
\]
But squares have Legendre value $1$, so
\[
1 = \left( \frac{D}{p} \right)\cdot 1.
\]
Thus
\[
\left( \frac{D}{p} \right) = 1,
\]
which means $D$ is a quadratic residue modulo $p$.

For the special case $p \mid (x^{2} - Dy^{2})$ and $(x,y)=1$, take $a=1$, $b=0$, $c=-D$ so
that $D=b^{2}-4ac=4D$. Since $4$ is a square modulo $p$, the conclusion that $D$ is a
square modulo $p$ follows immediately.

\end{proof}
\end{homeworkProblem}

\begin{homeworkProblem}
    We determine for which odd primes $p$ the quadratic congruence
    \[
        x^2 \equiv 2 \pmod{p}
    \]
    has a solution.  This holds exactly when
    \[
        \left( \frac{2}{p} \right) = 1.
    \]

    By Theorem 7.16 from the notes,
    \[
        \left( \frac{2}{p} \right)
        =
        \begin{cases}
            1 & \text{if } p \equiv 1 \text{ or } 7 \pmod{8}, \\[6pt]
            -1 & \text{if } p \equiv 3 \text{ or } 5 \pmod{8}.
        \end{cases}
    \]
    Therefore $2$ is a quadratic residue modulo $p$ precisely when
    \[
        p \equiv 1 \quad \text{or} \quad p \equiv 7 \pmod{8}.
    \]

\end{homeworkProblem}

\end{document}
